\subsection*{从圆周运动读出两个变化率}
单位圆上的位置写成 $(\cos\theta,\sin\theta)$。角度增加一点，点沿圆的切向移动；切向与半径垂直，相当于把位置方向转九十度，得到 $(-\sin\theta,\cos\theta)$。两坐标分别告诉我们横向与纵向变化速度。

在圆最右侧，纵向高度增长最快，正弦的变化率为一；来到最上方，高度到达峰值，一阶变化率为零；再向左上方走，高度开始下降。这样先预测符号，再看公式 $\cos\theta$，公式就变成了运动记录。

\subsection*{弧度不是可省略的单位约定}
为什么速度向量恰好有单位长度？因为单位圆上的弧度增量等于弧长增量。若用“度”作为输入，一个单位角度只走 $\pi/180$ 的弧长，导数必须多一个换算因子。导数不仅关乎函数形状，也关乎输入单位。

几何论证背后有一个极限：小弧段与切线方向的直线段越来越接近，例如 $\sin h/h\to1$。若用三角不等式或面积夹逼证明它，应在弧度约定下进行，不能把待证导数又拿来证明自身。

观看复平面动画时，同时观察位置箭头与垂直的速度箭头，再把它们分别投到横纵轴。你应能从图中解释正弦导数为何是余弦、余弦导数为何带负号。之后复合角度 $2x$ 的额外倍率，将由链式法则处理。
